% Resumo e palavras-chave no formato UFN
\begin{abstract}
A disciplina de Química no Ensino Médio, particularmente no que diz respeito à Tabela Periódica, enfrenta desafios críticos devido ao excesso de abstração e ao baixo engajamento dos estudantes. O objetivo principal deste projeto é desenvolver o Alchemon, um jogo sério que utiliza mecânicas de captura e colecionismo de criaturas para promover a aprendizagem tangencial e o engajamento volitivo dos estudantes. A base teórica baseia-se na fusão da Teoria da Autodeterminação, da Teoria da Aprendizagem Significativa e da Aprendizagem Tangencial, combinadas com a Taxonomia de Bloom para concentrar a carga cognitiva nos alunos. Os resultados deste trabalho demonstram a viabilidade técnica da proposta, validam o enquadramento conceitual e o cronograma necessários para o futuro desenvolvimento e validação prática do protótipo.
\end{abstract}

\textbf{\textit{Palavras-chave}: jogo sério; ensino de Química; aprendizagem tangencial; \textit{Game-Based Learning}; Tabela Periódica.}
